\par\medskip
\begin{center}
\textls[100]{d) Die nachgelassene Fortsetzung der Disquisitiones.}\\
\textls[100]{Funktionentheoretisches.}
\end{center}
\noindent\hspace*{\parindent}Am Schluss des Art. 3 der \textit{Disquisitiones} heisst es: »Weiter unten, vom
IV. Abschnitte ab, werden wir dagegen unsere Funktion einem höhern Prinzip
unterordnen, das eine ganz allgemeine Anwendung ermöglicht«. Dieses
Prinzip finden wir in dem ersten der von \textsc{Gauss} nicht veröffentlichten Artikel
(38, Werke III, S. 207), es ist die Differentialgleichung, die \textsc{Gauss} als besondern
Fall der Relationen zwischen verwandten Funktionen in der Form von
\textsc{Euler} (vergl. oben Gl. (5), S. 136)
\begingroup\makeatletter\tagsleft@true\let\veqno\leqno\makeatother
\begin{equation*}
0=\alpha\beta P-(\gamma-(\alpha+\beta+1)x)\frac{dP}{dx}
-(x-xx)\frac{d^2P}{dx^2}\tag{5\,a}
\end{equation*}\endgroup
\noindent ableitet. Diese sei — so sagt \textsc{Gauss} — als die exakte Definition seiner
Funktion anzusehen, der jedoch noch die für \(P=F(\alpha,\beta,\gamma,x)\) geltenden Anfangsbedingungen
\[
P=1,\quad \frac{dP}{dx}=\frac{\alpha\beta}{\gamma},\quad
\frac{d^2P}{dx^2}=\frac{\alpha\beta(\alpha+1)(\beta+1)}{\gamma(\gamma+1)}
\]
für \(x=0\) hinzuzufügen seien. »Auf diese Weise« — so fährt er fort — »wird
für jeden Wert von \(x\), zu dem man von \(x=0\) ausgehend durch kontinuierliche
Stufen übergeht, so jedoch, dass man den Wert \(x=1\), für den \(x-x^2=0\)
wird, nicht berührt, \(P\) eine völlig bestimmte Grösse sein; aber offenbar kann
man auf diese Weise zu reellen positiven Werten von \(x\), die grösser sind
als Eins, überhaupt nicht gelangen, es sei denn, dass man durch imaginäre
Werte hindurchgeht, und da dies auf unendlich verschiedene Weise geschehen
kann, ohne das Gebot der Kontinuität bei Seite zu lassen, so ist es möglich,\par
\clearpage\SourceLeaf{234}\ArithmeticHead{149}{R}{VIId. DIE FORTSETZUNG DER DISQUISITIONES.}
\noindent dass demselben Werte von \(x\) mehrere, ja sogar unendlich viele Werte von \(P\)
entsprechen, was ja auch für mehrere wohlbekannte Funktionen eintritt.«
Fast mit denselben Worten spricht sich \textsc{Gauss} in dem Briefe an \textsc{Bessel} vom
21. Nov. 1811 (Werke X, 1, S. 363, 364) aus, und fügt dort als Beispiel für
die »wohlbekannten Funktionen« noch den Logarithmus hinzu, dessen Definition
er mit \textsc{Acunha}\footnote{Wahrscheinlich ist gemeint \textsc{José Anastasio Cunha} (1744—1787) von dem posthum zwei Schriften: \textit{Principios mathematicos} (Lisbao 1790) und \textit{Ensaio sobre os principios de mechanica} (London 1807) erschienen sind.} durch Inversion der Reihe für die Exponentialfunktion
zu geben vorschlägt. Nach dieser prinzipiellen Bemerkung, auf die sich \textsc{Gauss}
vorbehält später ausführlich zurückzukommen, wird \(x\) wieder nicht grösser als
Eins vorausgesetzt und \(P\) weiter als die Summe der Reihe \(F(\alpha,\beta,\gamma,x)\) angesehen.
\par
\textsc{Gauss} bemerkt (Art. 39), dass die Differentialgleichung (5a) durch eine
lineare Transformation von \(x\) mit der allgemeineren Differentialgleichung (13),
wie sie oben S. 143 in dem Briefe an \textsc{Olbers} auftritt, identifiziert werden
kann. Die Transformation \(x=1-y\) liefert ein von \(F(\alpha,\beta,\gamma,x)\) linear unabhängiges
Integral
\[
F(\alpha,\beta,\alpha+\beta+1-\gamma,1-x),
\]
die \textsc{Eulersche} Substitution \(P=(1-x)^\mu P'\) (Art. 40) liefert für \(\mu=\gamma-\alpha-\beta\)
\[
F(\gamma-\alpha,\gamma-\beta,\gamma,x)=(1-x)^{\alpha+\beta-\gamma}F(\alpha,\beta,\gamma,x),
\]
die Substitution \(P=x^{1-\gamma}P'\) (Art. 41) endlich das Integral
\[
x^{1-\gamma}F(\alpha+1-\gamma,\beta+1-\gamma,2-\gamma,x),
\]
und nun wird die lineare Relation zwischen
\[
\begin{gathered}
F(\alpha,\beta,\gamma,x),\quad x^{1-\gamma}F(\alpha+1-\gamma,\beta+1-\gamma,2-\gamma,x),\\
F(\alpha,\beta,\alpha+\beta+1-\gamma,1-x)
\end{gathered}
\]
aufgestellt, und ihre Koeffizienten werden durch die \(\Pi\)-Funktion ausgedrückt
(Artikel 42—44).
\par
Die letztere Relation dient dazu, für Werte von \(x\) zwischen \(0{,}5\) und \(1\),
wo die Konvergenz der Reihe \(F(\alpha,\beta,\gamma,x)\) schwach ist, eine stärkere Konvergenz
zu erzielen, jedoch versagt dieses »Remedium« für ganzzahlige Werte
von \(\gamma\), wo jene Relation illusorisch wird (Art. 44). Diese Fälle werden dann
(Art. 45, 46) nach der Methode von \textsc{D’Alembert} erledigt; es tritt der Logarithmus\par
\clearpage\SourceLeaf{235}\ArithmeticHead{150}{L}{L. SCHLESINGER, ÜBER GAUSS' FUNKTIONENTHEORETISCHE ARBEITEN.}
\noindent auf, und als Beispiel wird die Reihe \(F(\tfrac12,\tfrac12,1,1-x)\) behandelt, die
(was \textsc{Gauss} a. a. O. nicht erwähnt) den reziproken Wert von \(M(1,\sqrt{x})\) darstellt.
Die Gleichung [90] Werke III, S. 217:
\[
\begin{aligned}
F\left(\frac12,\frac12,1,1-x\right)
=-\frac1\pi\Bigl\{&\log\frac1{16}x\mathbin{.}F\left(\frac12,\frac12,1,x\right)\\
&+\frac12x+\frac{21}{64}xx+\frac{185}{768}x^3+\ldots\Bigr\}
\end{aligned}
\]
ist ihrem Inhalte nach geradezu mit der Gleichung (7) Werke X, 1, S. 186,
der Scheda Ac (vergl. oben (10) Abschn. IVa), S. 68) identisch und ergibt
insbesondere für \(x=0\) die Gleichung
\[
\lim_{p\to0}M(1,p)\log\frac4p=\frac\pi2.
\]
\par
Es folgen (Art. 47) die Transformationen
\[
x=\frac{y}{y-1},\quad P=(1-y)^\mu P',\quad \mu=\alpha,\beta
\]
und (Art. 49)
\[
x=\frac1y,\quad P=y^\mu P',\quad \mu=\alpha,\beta,
\]
die weitere lineare Relationen zwischen \(F\)-Reihen liefern, dann (Art. 51, 52)
einige quadratische bezw. bilineare Relationen und endlich (Art. 53—55):
»Einige spezielle Theoreme« für spezielle Werte der \(\alpha,\beta,\gamma\), die auf die »paradoxe«
Gleichung (Werke III, S. 226)
\begingroup\renewcommand{\thefootnote}{216a)}
\[
F(2\alpha,2\beta,\alpha+\beta+\tfrac12,y)
=F(2\alpha,2\beta,\alpha+\beta+\tfrac12,1-y)\text{\footnotemark}
\]
\footnotetext{\textsc{Gauss} macht a. a. O. in der Reihe \(w_1=F(\alpha,\beta,\alpha+\beta+\tfrac12,x)\) die Substitution \(x=4y(1-y)\);
der Konvergenzbereich \(|4y(1-y)|<1\) besteht dann aus den beiden Schlingen der Lemniskate mit den
Brennpunkten \(0,1\) und dem Doppelpunkte \(\tfrac12\). In der \(y=0\) enthaltenden Schlinge stimmt \(w_1\) mit der Reihe
\(w_2=F(2\alpha,2\beta,\alpha+\beta+\tfrac12,y)\), in der \(y=1\) enthaltenden, mit der Reihe \(w_3=F(2\alpha,2\beta,\alpha+\beta+\tfrac12,1-y)\)
überein, die Reihe \(w_1=F(\alpha,\beta,\alpha+\beta+\tfrac12,4y-4y^2)\) stellt also in den beiden Schlingen zwei verschiedene
analytische Funktionen von \(y\) dar (vergl. J. \textsc{Thomae}, Göttinger Nachrichten, Math.-Phys. Klasse, 1904,
S. 465). Ein analytischer Zusammenhang zwischen \(w_2\) und \(w_3\) stellt sich her, wenn man diese Ausdrücke
als Funktionen von \(x\) betrachtet, d. h. (vergl. Werke III, S. 227) \(y=\tfrac12-\tfrac12\sqrt{1-x}\), also \(1-y=\tfrac12+\tfrac12\sqrt{1-x}\)
setzt; ein Umlauf um \(x=1\) führt dann \(w_2,w_3\) in einander über. Noch deutlicher wird dies, wenn man
(siehe Werke III, S. 227, Art. 56) gemäss der Formel (87), Werke III, S. 213, \(w_1\) durch die Gleichung
\[
\begin{aligned}
w_1={}&A\mathbin{.}F(\alpha,\beta,\tfrac12,1-x)\\
&+B(1-x)^{\frac12}F(\beta+\tfrac12,\alpha+\tfrac12,\tfrac32,1-x)
\end{aligned}
\]
in der Umgebung von \(x=1\) darstellt, wo \(A,B\) durch die \(\Pi\)-Funktion ausdrückbare Konstanten bedeuten.
Die in der gedachten Umgebung uniformisierende Substitution \(x=4y(1-y)\) lässt dann \(w_1\) in die beiden
getrennten, in der Umgebung von \(y=\tfrac12\) holomorphen analytischen Funktionen
\[
\begin{aligned}
w_2={}&A\mathbin{.}F(\alpha,\beta,\tfrac12,4(y-\tfrac12)^2)\\
&+2B(y-\tfrac12)\mathbin{.}F(\beta+\tfrac12,\alpha+\tfrac12,\tfrac32,4(y-\tfrac12)^2)\\
w_3={}&A\mathbin{.}F(\alpha,\beta,\tfrac12,4(y-\tfrac12)^2)\\
&-2B(y-\tfrac12)\mathbin{.}F(\beta+\tfrac12,\alpha+\tfrac12,\tfrac32,4(y-\tfrac12)^2)
\end{aligned}
\]
zerfallen.}
\addtocounter{footnote}{-1}\endgroup\par
\clearpage\SourceLeaf{236}\ArithmeticHead{151}{R}{VIId. DIE FORTSETZUNG DER DISQUISITIONES.}
\noindent führen. Um diese zu erklären, bemerkt \textsc{Gauss}: »Man muss bedenken, dass
zwischen zwei verschiedenen Bedeutungen der Charakteristik \(F\) unterschieden
werden muss, je nachdem diese nämlich entweder eine Funktion darstellt, deren
Eigenart durch die Differentialgleichung (5a) ausgedrückt wird, oder aber nur
die Summe der unendlichen Reihe. Die letztere gibt stets eine vollkommen
bestimmte Grösse, solange das vierte Element zwischen \(-1\) und \(+1\) gelegen
ist, \mbox{.\ .\ .\ .}\ die erstere Bedeutung dagegen stellt eine allgemeine Funktion dar,
die sich stets nach dem Gesetze der Kontinuität ändert, wenn das vierte Element
kontinuierlich geändert wird, gleichgültig ob ihm dabei reelle oder imaginäre
Werte beigelegt werden, wenn man nur allemal die Werte \(0\) und \(1\) vermeidet.
Hiernach ist klar, dass die Funktion in dem eben dargelegten Sinne
für gleiche Werte des vierten Elements, indem man einen Übergang oder
besser, einen Rückgang (reditus) durch imaginäre Werte vollzieht, verschiedene
Werte annehmen kann, unter denen einer derjenige ist, den die Reihe \(F\)
liefert.« Die paradoxe Gleichung ist also auf ähnliche Weise entstanden, »als
wollte man aus \(\arcsin\tfrac12=30^0\) und \(\arcsin\tfrac12=150^0\) schliessen \(30^0=150^0\)«.
\par
Wir sehen hieraus, dass \textsc{Gauss} zu jener Zeit über die Unendlichvieldeutigkeit
der Funktionen komplexer Variabeln, über ihre Wertänderung bei
Umläufen (er prägt selbst das Wort \textls[100]{reditus}) um singuläre Punkte vollständige
Klarheit hatte. Darüber hinaus hat er für die spezielle Differentialgleichung
(5a) — wie A. \textsc{Voss} sich ausdrückt — die Grundlagen derjenigen Theorie
der linearen Differentialgleichungen angedeutet, die allgemein durch die Arbeiten
von L. \textsc{Fuchs} begründet, sich das direkte Studium der Lösungen mit
Hilfe der Methoden der Funktionentheorie zur Aufgabe macht\footnote{A. \textsc{Voss}, \textit{Über das Wesen der Mathematik}, Leipzig 1908, S. 56. Dabei ist zu bemerken, dass der die nachgelassene Fortsetzung der \textit{Disquisitiones} enthaltende Bd. III der Werke erst nach dem Erscheinen der ersten Arbeiten von \textsc{Fuchs} (1865—68) veröffentlicht worden ist (vergl. Fussnote 5)).}.
\par
Für \textsc{Gauss'} damalige Kenntnisse in der Funktionentheorie kommt nun
noch die berühmte Stelle in dem Briefe an \textsc{Bessel} vom 18. Dezember 1811
(Werke X, 1, S. 366—368) in Betracht, wo von dem Integral \(\int\varphi(x)\,dx\)
für \(x=a+bi\) die Rede ist. Um dieses Integral zu definieren, lässt er \(x\)
durch unendlich kleine Inkremente von der Form \(\alpha+\beta i\) von dem Werte,
wo das Integral verschwindet, bis zu \(x=a+bi\) übergehn und summiert dann
über alle \(\varphi(x)\,dx\). Der stetige Übergang von einem Werte von \(x\) zu einem\par
\clearpage\SourceLeaf{237}\ArithmeticHead{152}{L}{L. SCHLESINGER, ÜBER GAUSS' FUNKTIONENTHEORETISCHE ARBEITEN.}
\noindent andern geschieht aber durch eine Linie in der Ebene, und ist mithin auf unendlich
viele Arten möglich. Es gilt dann der »sehr schöne Lehrsatz, \mbox{.\ .\ .\ .}
dass das Integral \(\int\varphi(x)\,dx\) nach zwei verschiedenen Übergängen immer einerlei
Wert erhalte, wenn innerhalb des zwischen beiden, die Übergänge repräsentierenden
Linien eingeschlossenen Flächenraumes nirgends \(\varphi(x)=\infty\) wird«
wobei noch vorausgesetzt wird, »dass \(\varphi(x)\) selbst eine einförmige Funktion von
\(x\) ist, oder wenigstens für deren Werte innerhalb jenes ganzen Flächenraumes
nur ein System von Werten, ohne Unterbrechung der Stetigkeit angenommen
wird«. Daraus ist auch klar, wie \(\int\varphi(x)\,dx\) für einen Wert von \(x\) »mehrere
Werte haben kann, indem man nämlich beim Übergange dahin um einen solchen
Punkt, wo \(\varphi(x)=\infty\) wird, entweder gar nicht oder einmal oder mehreremale herumgehen kann.« Als Beispiel nimmt er wieder den \(\log x=\int\frac{dx}{x}\) und
hält dagegen das mit dem Integrallogarithmus zusammenhängende Integral
\[
\int\frac{e^x-1}{x}\,dx,
\]
bei dem der Integrand für keinen endlichen Wert unendlich wird, so dass
das Integral eine »einförmige« Funktion ist, die ja auch durch die immer
konvergierende Potenzreihe \(x+\frac14x^2+\frac1{18}x^3+\ldots\) gegeben wird. »Ich wollte«
— fährt er fort — »Herr \textsc{Soldner} hätte \mbox{.\ .\ .}\ statt seines \(\operatorname{li}x=\int\frac{dx}{\log x}\) lieber
jene gewählt, da eine einförmige Funktion immer ohne Vergleich als klassischer
und einfacher anzusehen ist, als eine vielförmige, zumal da \(\log x\) selbst
schon eine vielförmige Funktion ist.«\par

